% !TeX root = ../main.tex

\documentclass[../main.tex]{subfiles}
\begin{document}

\chapter*{KẾT LUẬN VÀ HƯỚNG PHÁT TRIỂN}
\phantomsection
\addcontentsline{toc}{chapter}{KẾT LUẬN VÀ HƯỚNG PHÁT TRIỂN}
\label{chuong7}

\section*{Kết luận}

Đồ án đã hoàn thiện thuyết minh kỹ thuật cho phương án điều chỉnh công nghệ của Nhà máy xử lý chất thải rắn sinh hoạt thành phố Hội An, công suất 120 tấn/ngày đêm. Trên cơ sở đặc tính rác đầu vào, cân bằng vật chất -- năng lượng và yêu cầu vận hành thực tế, phương án chuyển từ khí hóa -- tuabin khí sang nhiệt phân -- chưng cất dầu -- phát điện diesel được đánh giá là phù hợp hơn với điều kiện của dự án.

Kết quả phân tích cho thấy CTRSH tại Hội An có tỷ lệ hữu cơ cao, độ ẩm trung bình 56,3\% và thành phần không đồng nhất. Các đặc điểm này làm giảm tính phù hợp của công nghệ khí hóa, trong khi công nghệ nhiệt phân cho phép chuyển hóa phần RDF và rác nhựa sau tiền xử lý thành dầu nhiệt phân, khí không ngưng và than carbon -- trong đó dầu lỏng có ưu điểm là có thể lưu trữ, chưng cất và sử dụng linh hoạt hơn khí tổng hợp.

Về cân bằng vật chất, từ 120 tấn CTRSH/ngày.đêm, khối tiền xử lý tạo ra khoảng 46,08 tấn RDF/ngày cấp lò nhiệt phân. Kết quả thiết kế sơ bộ buồng nhiệt phân quay (Mục~\ref{sec:tinh-toan-lo-nhiet-phan-quay}) cho thấy một mô-đun lò có đường kính trong $D_i = 2{,}0$ m, chiều dài $L = 8{,}0$ m, thể tích khoảng 25 m$^3$ và thời gian lưu thiết kế 6 giờ, với tổng nhiệt cần cấp khoảng 489 kW (Bảng~\ref{tab:ket-qua-thiet-ke-lo}). Dòng nhiệt phân dự kiến gồm 20,74 tấn/ngày dầu thô, 13,82 tấn/ngày khí không ngưng và 11,52 tấn/ngày than carbon.

Về cân bằng năng lượng, công suất điện tinh đạt khoảng 1,82--1,84 MW, tương ứng sản lượng điện ròng khoảng 15.943 MWh/năm, khẳng định nhà máy có khả năng thu hồi năng lượng ở quy mô công nghiệp. Phân tích độ nhạy xác nhận phương án vẫn khả thi ngay ở điều kiện bất lợi nhất.

Về hiệu quả kinh tế, tổng mức đầu tư CAPEX khái toán khoảng 214,8 tỷ đồng (tương đương 8,6 triệu USD), với thời gian hoàn vốn khoảng 5,3 năm, NPV tại r = 8\% đạt khoảng 185 tỷ đồng và IRR khoảng 13,5\% vượt chi phí vốn. Các chỉ số tài chính này chứng minh dự án khả thi về mặt tài chính.

Phương án nhiệt phân kết hợp chưng cất dầu và phát điện diesel có tính khả thi kỹ thuật và kinh tế đối với điều kiện rác sinh hoạt Hội An, đồng thời phù hợp với định hướng giảm chôn lấp và tận dụng năng lượng từ chất thải theo chính sách nhà nước hiện hành.

\section*{Kiến nghị}

Để phương án công nghệ đạt hiệu quả trong thực tế, cần chú trọng ba nhóm vấn đề chính. Nhóm thứ nhất là nguyên liệu đầu vào: duy trì chất lượng RDF ổn định thông qua thúc đẩy phân loại rác tại nguồn, ký hợp đồng cung ứng dài hạn và xây dựng khu đệm RDF đủ cho ít nhất 2 ngày vận hành. Nhóm thứ hai là vận hành và an toàn: ban hành đầy đủ SOP cho từng công đoạn, đào tạo nhân lực trước khi vận hành thương mại, kiểm tra định kỳ hệ thống ESD, cảm biến LEL và phương tiện PCCC. Nhóm thứ ba là kiểm chứng thực tế: trong giai đoạn chạy thử, cần thu thập số liệu thực tế về thành phần rác, độ ẩm RDF, chất lượng dầu, tiêu hao năng lượng và hiệu quả xử lý môi trường để hiệu chỉnh thiết kế và hoàn thiện hồ sơ kỹ thuật.

\section*{Hướng phát triển tiếp theo}

Trong quá trình triển khai thực tế và nghiên cứu tiếp theo, các hướng phát triển tiềm năng bao gồm:

\begin{itemize}
    \item \textbf{Nghiên cứu tối ưu hóa quá trình nhiệt phân:} Thử nghiệm các điều kiện vận hành khác nhau (nhiệt độ, thời gian lưu, tốc độ gia nhiệt) để tối đa hóa tỷ lệ dầu và chất lượng sản phẩm cho từng loại nguyên liệu RDF cụ thể tại Hội An.
    \item \textbf{Nghiên cứu nâng cao chất lượng dầu nhiệt phân:} Áp dụng các phương pháp tinh chế tiên tiến (hydrotreating, catalytic cracking) để cải thiện trị số cetane và giảm hàm lượng lưu huỳnh, cho phép tăng tỷ lệ pha trộn dầu nhiệt phân trong nhiên liệu diesel.
    \item \textbf{Tích hợp với hệ thống năng lượng mặt trời:} Kết hợp phát điện từ nhiệt phân với hệ thống điện mặt trời áp mái để tăng tổng công suất phát điện và giảm phụ thuộc vào lưới điện.
    \item \textbf{Nghiên cứu thu hồi kim loại từ tro và than carbon:} Khảo sát khả năng thu hồi các kim loại có giá trị (nhôm, đồng, vàng) từ tro bay và tro đáy, cũng như nghiên cứu ứng dụng than carbon từ nhiệt phân làm vật liệu hấp phụ hoặc nguyên liệu cho các ngành công nghiệp khác.
\end{itemize}

\end{document}
